%  LaTeX support: latex@mdpi.com 
%  For support, please attach all files needed for compiling as well as the log file, and specify your operating system, LaTeX version, and LaTeX editor.

%=================================================================
\documentclass[notspecified,article,submit,moreauthors]{Definitions/mdpi} 
%\documentclass[journal,article,submit,moreauthors,unicode]{Definitions/mdpi} % For some letters Palatino lacks
%\documentclass[preprints,article,submit,moreauthors]{Definitions/mdpi} 
%\documentclass[preprints,article,submit,moreauthors,unicode]{Definitions/mdpi} 
% For posting an early version of this manuscript as a preprint, you may use "preprints" as the journal. Changing "submit" to "accept" before posting will remove line numbers.

%--------------------
% Class Options:
%--------------------
%----------
% journal
%----------
% Choose between the following MDPI journals:
% accountaudit, acn, acoustics, actuators, addictions, addictprev, adhesives, adolescents, aee, aerobiology, aeronautics, aerospace, agriculture, agriengineering, agrochemicals, agronomy, ai, aiagents, aichem, aieng, aimater, aimed, aipa, air, aisens, aisoc, algorithms, allergies, alloys, amh, analog, analytica, analytics, anatomia, anesthres, animals, antibiotics, antibodies, antioxidants, applbiosci, appliedchem, appliedmath, appliedphys, applmech, applmicrobiol, applnano, applsci, aps, aquacj, archaeolstud, architecture, arm, arthropoda, asc, asd, asi, astronautics, astronomy, atmosphere, atoms, audiolres, automation, axioms, bacteria, batteries, bcrc, bdcc, beverages, bioactives, biochem, bioengineering, biologics, biology, biomass, biomechanics, biomed, biomedicines, biomedinformatics, biomimetics, biomolecules, biophysica, bioresourbioprod, biosensors, biosphere, biotech, birds, blockchains, bloods, blsf, brainsci, breath, buildings, cancers, carbon, cardio, cardiogenetics, cardiovascmed, catalysts, cells, ceramics, challenges, chemeco, chemengineering, chemistry, chemosensors, chemproc, children, chips, cimb, civileng, cleantechnol, climate, clinbioenerg, clinpract, clockssleep, cmd, cmtr, coasts, coatings, colloids, colorants, commodities, complexities, complications, compounds, computation, computers, condensedmatter, conservation, constrmater, cosmetics, covid, crops, cryo, cryptography, crystals, csmf, ctn, culture, curroncol, cyber, dairy, data, ddc, dentistry, dermato, dermatopathology, designs, devices, dgg, dhi, diabetology, diagnostics, dietetics, digital, disabilities, diseases, diversity, dna, drones, drugreact, drylands, dynamics, earth, ebj, ece, ecm, ecologies, edm, eesp, electricity, electrochem, electronicmat, electronics, embodiedintell, encyclopedia, endocrines, energies, eng, engproc, entomology, entropy, environments, environremediat, enzymology, epidemiologia, epigenomes, epilepsy, esa, est, fermentation, fibers, fintech, fire, fishes, fluids, foodeng, foods, forecasting, forensicsci, forests, fossstud, foundations, fractalfract, freshwater, fuels, future, futureinternet, futurepharmacol, futurephys, futuretransp, galaxies, gases, gastroent, gastrointestdisord, gastronomy, gels, genes, geochem, geographies, geohazards, geomatics, geometry, geosciences, geotechnics, geriatrics, germs, glacies, grainsci, grasses, green, greenhealth, gucdd, gynecolcancer, hardware, hazardousmatters, healthcare, healthcmater, hearts, hemato, hematolrep, hep, heritage, higheredu, highthroughput, horticulturae, hospitals, humanoids, hydrobiology, hydrogen, hydrology, hydrometeorology, hydropower, hygiene, idr, iic, ijem, ijerph, ijgi, ijmd, ijms, ijns, ijom, ijp, ijpb, ijt, ijtm, ijtpp, ijtst, ime, immuno, industries, inflammj, informatics, information, infrastructures, inorganics, insects, instruments, inventions, iot, iron, j, jaestheticmed, jal, japma, jcdd, jcm, jcp, jcrm, jcs, jcto, jdad, jdb, jdcp, jdream, jemr, jeta, jfb, jfmk, jfth, jgbg, jgg, jimaging, jlpea, jmahp, jmmp, jmms, jmp, jmse, jne, jnt, jof, joitmc, joma, jop, joptm, jor, jox, jpbi, jphytomed, jpm, jsan, jtaer, jvd, jzbg, kidneydial, kinasesphosphatases, knowledge, labmed, laboratories, lam, land, langtechnol, lasers, legumes, life, lights, limnolrev, lipidology, liquids, livers, logics, logistics, lubricants, lymphatics, machines, macromol, magnetism, magnetochemistry, make, marinedrugs, maritime, materials, materproc, mathematics, mca, mcb, measurements, meb, medicina, medicines, medsci, membranes, merits, metabolites, metals, metamaterialsadv, meteorology, methane, metrics, metrology, micro, microarrays, microbiolres, microelectronics, micromachines, microorganisms, microplastics, microwave, minerals, mining, mmphys, modelling, molbank, molecules, mollusca, mps, msf, mti, multimedia, muscles, mutation, nanoenergyadv, nanomanufacturing, nanomaterials, ncrna, ndt, network, neuroglia, neuroimaging, neurolint, neurosci, nitrogen, notspecified, nursrep, nutraceuticals, nutrients, obesities, occuphealth, oceans, ohbm, onco, optics, optoelectronics, oral, organics, organoids, osteology, oxygen, pandemics, parasites, parasitologia, particles, pathogens, pathophysiology, pediatrrep, pets, pharmaceuticals, pharmaceutics, pharmacoepidemiology, pharmacy, phc, philosophies, photochem, photonics, photovoltaics, phycology, physchem, physics, physiologia, plants, plasma, platforms, pnr, pollutants, polymers, polysaccharides, populations, poultry, powders, precision, precisoncol, preprints, prevhealthc, proceedings, processes, prosthesis, proteomes, psf, psychiatryint, psychoactives, purification, quantumrep, quaternary, qubs, radiation, rareearths, rdt, reactions, realestate, receptors, recycling, regeneration, remotesensing, reports, reprodmed, resources, rheumato, rjpm, robotics, rsee, ruminants, safety, sanpp, sca, sci, scipharm, sclerosis, seeds, sensors, senssyst, separations, sexes, shi, signals, sinusitis, siuj, skins, smartcities, smartfish, smartgrids, sna, societies, software, soilsystems, solar, solids, spectroscj, sports, standards, stats, std, stratsediment, stresses, strokerep, superintelligence, surfaces, surgeries, susagric, suschem, sustainability, sustainfoods, sustainpackag, symmetry, synbio, systems, tae, targets, taxonomy, technologies, telecom, test, textiles, thalassrep, therapeutics, thermo, timespace, tomography, toxics, toxins, tph, transplantology, transportation, traumacare, traumas, tri, tropicalmed, universe, urbansci, uro, vaccines, vehicles, venereology, vetsci, vibration, virtualworlds, viruses, vision, waste, water, welding, wem, wevj, wild, wind, women, world, zoonoticdis

%---------
% article
%---------
% The default type of manuscript is "article", but can be replaced by: 
% abstract, addendum, article, benchmark, book, bookreview, briefcommunication, briefreport, casereport, changes, clinicopathologicalchallenge, comment, commentary, communication, conceptpaper, conferenceproceedings, correction, conferencereport, creative, datadescriptor, discussion, entry, expressionofconcern, extendedabstract, editorial, essay, erratum, fieldguide, hypothesis, interestingimages, letter, meetingreport, monograph, newbookreceived, obituary, opinion, proceedingpaper, projectreport, reply, retraction, review, perspective, protocol, shortnote, studyprotocol, supfile, systematicreview, technicalnote, viewpoint, guidelines, registeredreport, tutorial, giantsinurology, urologyaroundtheworld, patentsummary

% supfile = supplementary materials

%----------
% submit
%----------
% The class option "submit" will be changed to "accept" by the Editorial Office when the paper is accepted. This will only make changes to the frontpage (e.g., the logo of the journal will get visible), the headings, and the copyright information. Also, line numbering will be removed. Journal info and pagination for accepted papers will also be assigned by the Editorial Office.

%------------------
% moreauthors
%------------------
% If there is only one author the class option oneauthor should be used. Otherwise use the class option moreauthors.

%=================================================================
% MDPI internal commands - do not modify
\firstpage{1} 
\makeatletter 
\setcounter{page}{\@firstpage} 
\makeatother
\pubvolume{1}
\issuenum{1}
\articlenumber{0}
\pubyear{2026}
\copyrightyear{2026}
%\externaleditor{Firstname Lastname} % More than 1 editor, please add `` and '' before the last editor name
\datereceived{ } 
\daterevised{ } % Comment out if no revised date
\dateaccepted{ } 
\datepublished{ } 
%\datecorrected{} % For corrected papers: "Corrected: XXX" date in the original paper.
%\dateretracted{} % For retracted papers: "Retracted: XXX" date in the original paper.
%\doinum{} % Used for some special journals, like molbank
%\pdfoutput=1 % Uncommented for upload to arXiv.org
%\CorrStatement{yes}  % For updates
%\longauthorlist{yes} % For many authors that exceed the left citation part
%\IsAssociation{yes} % For association journals

%=================================================================
% Add packages and commands here. The following packages are loaded in our class file: fontenc, inputenc, calc, indentfirst, fancyhdr, graphicx, epstopdf, lastpage, ifthen, float, amsmath, amssymb, lineno, setspace, enumitem, mathpazo, booktabs, titlesec, etoolbox, tabto, xcolor, colortbl, soul, multirow, microtype, tikz, totcount, changepage, attrib, upgreek, array, tabularx, pbox, ragged2e, tocloft, marginnote, marginfix, enotez, amsthm, natbib, hyperref, cleveref, scrextend, url, geometry, newfloat, caption, draftwatermark, seqsplit
% cleveref: load \crefname definitions after \begin{document}

%=================================================================
% Please use the following mathematics environments: Theorem, Lemma, Corollary, Proposition, Characterization, Property, Problem, Example, ExamplesandDefinitions, Hypothesis, Remark, Definition, Notation, Assumption
%% For proofs, please use the proof environment (the amsthm package is loaded by the MDPI class).

%=================================================================

% Additional package required by the migrated source.

\usepackage{placeins}
\setlength{\emergencystretch}{2em}

%=================================================================
% Manuscript metadata migrated from the source paper.
\Title{Auditing Hospital-Associated Context Sensitivity with a
Pixel-Identical Label-Defining Region}

% Replace the placeholders below with the final author details before submission.
\Author{Given Name Surname $^{1}$, Given Name Surname $^{2}$, Given Name Surname $^{3}$, Given Name Surname $^{4}$, Given Name Surname $^{5}$, Given Name Surname $^{6}$}
\AuthorNames{Given Name Surname, Given Name Surname, Given Name Surname, Given Name Surname, Given Name Surname and Given Name Surname}
\address{%
$^{1}$ \quad dept. name of organization (of Aff.), name of organization (of Aff.), City, Country; email address or ORCID\\
$^{2}$ \quad dept. name of organization (of Aff.), name of organization (of Aff.), City, Country; email address or ORCID\\
$^{3}$ \quad dept. name of organization (of Aff.), name of organization (of Aff.), City, Country; email address or ORCID\\
$^{4}$ \quad dept. name of organization (of Aff.), name of organization (of Aff.), City, Country; email address or ORCID\\
$^{5}$ \quad dept. name of organization (of Aff.), name of organization (of Aff.), City, Country; email address or ORCID\\
$^{6}$ \quad dept. name of organization (of Aff.), name of organization (of Aff.), City, Country; email address or ORCID}
\corres{Correspondence: To be completed by the authors.}

\abstract{Auditing hospital-associated context use is difficult because
external validation measures distribution shift without isolating the evidence
used by a fixed model. Existing attribution and domain-analysis methods provide evidence, but
they do not separate hospital-indexed context sensitivity from generic image
replacement effects. To address this limitation, we propose MedStyleAudit, a
paired counterfactual audit that keeps the central 32-by-32-pixel
label-defining region of each Camelyon17-WILDS patch pixel-identical while
replacing its periphery with matched tissue from the same or another hospital.
The framework combines balanced triplet matching, a same-hospital control,
Hospital Context Sensitivity (HCS), Hospital Context Effect (HCE), and
prespecified coverage and construction checks. On hospital-1 validation, HCS
was negative for all three ResNet-50 seeds (mean -0.0906, SD 0.0097), with every
95\% interval excluding zero: cross-hospital context was less destabilizing
than matched same-hospital replacement. The result persisted with an 8-pixel
buffer and hard compositing, while region-only controls were exactly zero. The
frozen matcher did not meet hospital-2 coverage; its post-hoc estimates were
positive but seed-heterogeneous (mean 0.1180, SD 0.0694). These results show
that hospital-indexed context sensitivity is population-dependent and that
MedStyleAudit should be interpreted as a scoped model audit rather than causal
identification of scanner, stain, or hospital effects.}
\keyword{medical image analysis; hospital-associated context; context sensitivity; counterfactual auditing; histopathology; domain shift}

\begin{document}

\begin{figure}[H]
\centering
\includegraphics[width=\columnwidth,height=4.2cm,keepaspectratio]{figures/figure1_motivation.png}
\caption{Motivation for a target-preserving context audit. A source patch is
paired with same-hospital (blue) and cross-hospital (orange) peripheral
contexts while the label-defining central ROI remains fixed. Both composites
are evaluated by the same classifier, enabling a paired comparison of context
sensitivity.}
\label{fig:motivation}
\end{figure}

\section{Introduction}

Medical image datasets are shaped by where specimens are collected, how they
are prepared, which scanners digitize them, and which patients enter the
cohort. These processes leave hospital-associated signals in radiographs and
histopathology
\cite{zech2018variable,badgeley2019confounding,degrave2021shortcuts,
howard2021site}. When those signals correlate with labels, a classifier can
achieve high internal accuracy while relying on evidence that is unstable or
clinically unintended~\cite{geirhos2020shortcut,koh2021wilds}.
More broadly, medical-imaging models can encode patient attributes that are not
readily apparent to human readers, reinforcing the need to distinguish
information availability from its use in a target decision
\cite{gichoya2022race}.

Existing studies examine this problem through external validation,
hospital-prediction probes, and attribution maps. These approaches establish
that performance changes across sites or that site information is present in
an image representation. However, a performance gap can also reflect case mix,
prevalence, annotation, or selection, and encoded site information need not
affect the target decision. Consequently, observational shift evidence does not
show whether a fixed classifier functionally uses a specified image context.

Direct testing requires a target-preserving intervention, but histopathology
makes such interventions difficult. Replacing peripheral tissue can introduce
a seam, change tissue composition, or select morphologically different donors.
An original-versus-cross-hospital comparison would therefore mix the response
to hospital-indexed context with the response to image replacement itself. A
valid audit must preserve the label-defining evidence, apply the same
construction to both comparison arms, and balance the content supplied by the
donors.

Building on this requirement, we introduce MedStyleAudit, the paired workflow
illustrated in Fig.~\ref{fig:motivation}. Camelyon17-WILDS labels each
$96\times96$ patch from its central $32\times32$ region
\cite{bandi2018camelyon17,koh2021wilds}. MedStyleAudit keeps that region
pixel-identical, replaces only the periphery, and pairs every cross-hospital
transplant with a content-matched same-hospital transplant. HCS measures the
cross-minus-within difference in absolute prediction change, while HCE retains
the signed direction. Coverage, balance, buffer, boundary, and ROI-only checks
define when these measures are interpretable.

The evaluation produces a result that a cross-hospital-only analysis would
miss. On hospital-1 validation, every model seed and donor-hospital direction
has negative HCS with a 95\% interval excluding zero: matched same-hospital
replacement is more destabilizing than cross-hospital replacement. The frozen
matcher fails its hospital-2 coverage gate, and the adapted post-hoc analysis
produces positive but seed-heterogeneous estimates. Hospital-indexed context
sensitivity is therefore not uniform across the two source populations.

\begin{samepage}
Following the MedStyleAudit design, the main contributions of this paper are
summarized as follows:
\begin{itemize}
    \item We formulate hospital-associated context use as a paired,
    target-preserving model audit and define HCS and HCE to separate excess
    instability from its signed direction.
    \item We develop Balanced Matched-Triplet construction with a
    same-hospital replacement control and explicit coverage, balance,
    pixel-identity, buffer, and boundary checks.
    \item We evaluate three ResNet-50 seeds across confirmatory and exploratory
    hospital populations, showing a stable negative H1 response and a positive
    but heterogeneous post-hoc H2 response.
\end{itemize}
\end{samepage}

\section{Related Work}

\subsection{Medical Shortcuts and Hospital-Indexed Variation}

Shortcut learning describes reliance on predictive but semantically
unintended features~\cite{geirhos2020shortcut}. Healthcare process and
source variables have confounded radiographic models
\cite{zech2018variable,badgeley2019confounding,degrave2021shortcuts}, and
direct shortcut tests assess whether a model uses a specified attribute
\cite{brown2023shortcut}. In digital pathology, tissue preparation,
staining, scanning, and cohort composition create site signatures that
affect representations, accuracy, and measured bias
\cite{howard2021site,stacke2021domain,bidgoli2022bias,tellez2019stain,
aubreville2023midog,komen2026robust}.
Hospital predictability, however, establishes information availability;
it does not show use by a tumor classifier.

This distinction has become central to shortcut research. A nuisance
attribute can be predictable from pixels without materially changing a target
decision, while a weakly predictable attribute may still influence a narrow
subset of clinically important examples. Correlation analyses and domain
classifiers therefore answer a representation question---whether site
information is encoded---rather than an intervention question---whether the
target predictor functionally depends on that information. Direct shortcut
tests narrow the gap by specifying an attribute and measuring its use
\cite{brown2023shortcut}, but their validity still depends on whether the
intervention preserves the target-relevant evidence.

Histopathology makes this separation difficult. Staining, scanner response,
tissue processing, morphology, and disease composition interact at several
spatial scales. Site effects can be global color transformations, local
texture changes, or differences in which tissue enters the field of view
\cite{stacke2021domain,howard2021site,bidgoli2022bias}. A method that removes
all hospital-predictive variation may consequently remove genuine pathology;
a method that changes appearance alone may miss hospital-associated
composition. Our audit addresses a narrower question by defining the
intervention spatially and balancing measured tissue geometry rather than
claiming complete removal of site information.

Interpretability tools can reveal recurring model strategies and expose
artifact-driven behavior, as illustrated by analyses of ``Clever Hans''
predictors~\cite{lapuschkin2019cleverhans}. Nevertheless, an attribution map
does not by itself establish that the highlighted signal is necessary for a
particular prediction, and this limitation is especially consequential in
health care~\cite{ghassemi2021falsehope}. These observations motivate an audit
based on controlled prediction changes rather than visual explanation alone.

\subsection{Domain Generalization Versus Model Auditing}

Domain-generalization benchmarks quantify how training procedures transfer
to held-out environments. WILDS formalizes this setting through distribution
shifts with explicit domain metadata~\cite{koh2021wilds}, and pathology work
has examined robustness to site, stain, and scanner variation
\cite{tellez2019stain,stacke2021domain,aubreville2023midog,
komen2026robust}. Such evaluations are indispensable
for deployment, but their primary estimand is predictive performance under a
new distribution. They do not isolate the image region or attribute that
causes a particular model response.

Model auditing instead holds the trained predictor fixed and probes a
specified behavior. The objective is explanatory and diagnostic rather than
optimization-oriented: the audit asks whether a controlled change in
candidate context alters the prediction. This difference also changes what
counts as a valid negative result. In domain generalization, a negative result
usually means failure to improve held-out accuracy. In an audit, a negative
or sign-reversed contrast can falsify the proposed shortcut mechanism while
still revealing substantial sensitivity to the control intervention.

\begin{figure}[H]
\centering
\includegraphics[width=\columnwidth,height=3.8cm,keepaspectratio]{figures/figure2_related_work_taxonomy.png}
\caption{Conceptual placement of MedStyleAudit. From left to right, the
taxonomy distinguishes observational domain-shift evidence, representational
site probes, attribution-based localization, and controlled functional-use
tests. The highlighted final module denotes the target-preserving
intervention studied here.}
\label{fig:related_taxonomy}
\end{figure}

\subsection{Counterfactual Auditing in Pathology}

Medical counterfactual methods generate decision-changing images or add
and remove candidate shortcuts~\cite{singla2023counterfactual,
weng2024fast}. HIPPO modifies tissue-patch sets in whole-slide models to
support quantitative hypothesis tests, including planted-bias analyses
\cite{kaczmarzyk2024hippo}. MoPaDi uses diffusion models to alter
morphology or stain-associated features and explain pathology classifiers
\cite{zigutyte2026mopadi}. Domain counterfactuals and
counterfactual-contrastive learning address related domain changes under
generative or latent-model assumptions
\cite{zhou2024domaincounterfactuals,roschewitz2025counterfactual}.

Controlled foreground--background studies outside medicine further show that
high-accuracy classifiers may respond strongly to background content and that
this behavior varies across architectures and training procedures
\cite{xiao2021backgrounds,moayeri2022sensitivity}. Although pathology context
is not equivalent to a natural-image background, these studies support the
general design principle of perturbing candidate context while preserving the
label-defining foreground.

Generative counterfactuals offer expressive interventions, but they introduce
an additional model whose fidelity and disentanglement must be assessed.
Generated changes may alter task evidence together with the intended factor,
and a visually plausible sample need not be causally faithful. Patch-removal
or replacement methods avoid some generative assumptions, yet a target model
can respond to the removal artifact itself. For this reason, a suitable
baseline must undergo the same construction process as the hypothesized
shortcut intervention.

Our paired design follows that principle. Both within- and cross-hospital
arms replace the same peripheral coordinates, use the same source ROI,
preserve label exactly, and apply the same feathering or hard boundary. The
arms differ in the hospital membership of the donor pool after balancing
prespecified descriptors. Subtracting the within-hospital arm therefore
controls generic replacement response more directly than an original-versus-
cross comparison, although residual unmeasured pathology and acquisition
differences remain.

\subsection{Governance of Public Medical-Imaging Data}
\label{subsec:ethics_related}

Public availability does not remove the ethical obligations attached to
human tissue and clinical metadata. The original CAMELYON collection was
approved by the local ethics committee of Radboud University Medical Center
(Commissie Mensgebonden Onderzoek regio Arnhem--Nijmegen; protocol
2016-2761), and informed consent was waived
\cite{litjens2018camelyon}. WILDS subsequently released a de-identified,
patch-based Camelyon17 variant with hospital and slide metadata under a CC0
public-domain dedication~\cite{koh2021wilds}.

De-identification and an open license do not exhaust the obligations of
secondary medical-data use. Medical big-data governance also concerns consent,
equity, downstream discrimination, and patient control
\cite{price2019privacy,vayena2018ethics}. Privacy-preserving and federated
learning can reduce some risks when data are distributed across institutions,
but technical safeguards complement rather than replace ethical and
institutional oversight~\cite{kaissis2020privacy}.

This study uses only the released images and metadata, performs no new
participant recruitment or specimen collection, and makes no attempt to
re-identify individuals or link records to external sources. We report
aggregate model behavior rather than patient-level findings. These safeguards
reduce privacy risk but do not replace institutional oversight. Before
submission, the authors must document whether their institution classified
this secondary analysis as approved, exempt, or outside the scope of
human-subjects review. The corresponding determination belongs in the
Institutional Review Board Statement.

\subsection{Positioning and Scope}

MedStyleAudit occupies the target-preserving intervention branch of
Fig.~\ref{fig:related_taxonomy}. It does not synthesize a new stain, infer a
latent acquisition factor, or retrain the classifier. Instead, it exchanges
observed peripheral tissue while keeping benchmark-defining evidence
pixel-exact and pairs each cross-hospital intervention with a comparably
matched same-hospital replacement. Matching makes the comparison conditional
on measured geometry, and coverage explicitly reports where that conditional
comparison is supported.

The resulting quantities are sensitivity contrasts rather than causal
effects of a hospital, scanner, or stain. This limited scope is deliberate.
It permits a reproducible functional-use test without asserting that hospital
identity is a manipulable mechanism or that the transplanted image represents
a physically realizable rescan. The work is therefore complementary to
domain-generalization training, representation analysis, and generative
counterfactual explanation rather than a replacement for any of them.

\section{Materials and Methods}
\label{sec:methods}

\subsection{Study Design and Dataset}

We conducted a paired counterfactual audit of three frozen tumor classifiers.
For each source patch, we preserved the central label-defining region and
replaced only the peripheral context. A same-hospital donor provided the
replacement control, and a content-matched donor from another hospital
provided the comparison arm. Hospital, slide, and patient metadata were used
only for splitting, matching, and grouped analysis; they were never classifier
inputs. Accordingly, all reported effects are hospital-indexed sensitivity
contrasts rather than causal effects of hospitals, scanners, or stains.

Camelyon17-WILDS contains 455,954 labeled $96\times96$ histopathology patches
from 50 whole-slide images across five hospitals
\cite{bandi2018camelyon17,koh2021wilds,litjens2018camelyon}. A patch is positive
when tumor occurs in its central $32\times32$ region. Table~\ref{tab:dataset_splits}
summarizes the released zero-indexed split metadata; the numbering displayed in
the WILDS paper differs from these API identifiers.

\begin{table}[H]
\caption{Camelyon17-WILDS split and domain information. Counts refer to the
labeled WILDS release.}
\label{tab:dataset_splits}
\centering
\scriptsize
\setlength{\tabcolsep}{3.2pt}
\begin{tabularx}{\columnwidth}{@{}lccrX@{}}
\toprule
Split & Key & Hospitals & Patches & Role \\
\midrule
Training & \texttt{train} & $\{0,3,4\}$ & 302,436 & training; cross-donor bank \\
ID validation & \texttt{id\_val} & $\{0,3,4\}$ & 33,560 & model selection; logit scale \\
OOD validation & \texttt{val} & $\{1\}$ & 34,904 & confirmatory H1 audit \\
OOD test & \texttt{test} & $\{2\}$ & 85,054 & exploratory H2 audit \\
\midrule
\multicolumn{3}{r}{Total} & 455,954 & 50 whole-slide images \\
\bottomrule
\end{tabularx}
\end{table}

The original CAMELYON collection was conducted under local ethics approval,
as documented by Litjens et al.~\cite{litjens2018camelyon}. The present study
used only the public, de-identified benchmark and did not recruit participants,
collect specimens, or access direct identifiers. Dataset licensing, secondary-use
oversight, and the required manuscript declarations are addressed in
Section~\ref{subsec:ethics_related} and the back matter.

\subsection{Classifier and Frozen Protocol}

We trained three ResNet-50 classifiers from scratch with seeds 11, 42, and 101
\cite{he2016resnet}. All models used native $96\times96$ inputs and the same
optimizer, schedule, augmentation, and H1 OOD-validation AUROC model-selection
rule. Best-checkpoint H1 OOD-validation AUROC was 0.8688, 0.8832, and 0.8983,
respectively ($0.8834\pm0.0148$). These model-selection scores are distinct
from the ID-validation and H2 OOD-test AUROCs. The descriptor, matching
parameters, audit subset, metrics, and training protocol were frozen before
hospital-2 predictions were accessed. Supplementary
Table~\ref{tab:full_implementation} lists the complete implementation settings.

\subsection{Pixel-Exact ROI Preservation}
\label{subsec:pixel_preservation}

Let $\Omega_R$ denote the central $32\times32$ ROI and $P_r$ the ROI plus a
square source buffer of radius $r$. A cosine feather of width $w$ defines a
source-pixel weight $\boldsymbol{\alpha}_{r,w}$ that equals one on $P_r$ and
decreases to zero outside the transition ring. For source image $\mathbf{x}_i$
and donor image $\mathbf{x}_j$, the composite is
\begin{equation}
\widetilde{\mathbf{x}}_{i\leftarrow j}^{(r,w)}
=\boldsymbol{\alpha}_{r,w}\odot\mathbf{x}_i
+(1-\boldsymbol{\alpha}_{r,w})\odot\mathbf{x}_j,
\qquad
\widetilde{\mathbf{x}}_{i\leftarrow j}^{(r,w)}[\Omega_R]
=\mathbf{x}_i[\Omega_R].
\label{eq:pixel_identity}
\end{equation}
The equality was asserted on the final classifier tensor after deterministic
preprocessing. The primary analysis used $r=0$ and $w=4$; robustness analyses
used an 8-pixel source buffer and a hard boundary while retaining the same
sources and donor triplets. Figure~\ref{fig:roi_construction} summarizes the
pixel-preserving construction.

\begin{figure}[H]
\centering
\includegraphics[width=\columnwidth,height=4.0cm,keepaspectratio]{figures/figure3_roi_construction.png}
\caption{Counterfactual construction. The central ROI and optional buffer
retain source pixels. The same blending rule is applied to matched
same-hospital and cross-hospital donors.}
\label{fig:roi_construction}
\end{figure}

\subsection{Content Matching and Donor Selection}
\label{subsec:content_descriptor}
\label{subsec:donor_selection}

We matched tissue geometry rather than color or texture. A fixed HSV-based
tissue mask used a $3\times3$ median filter, saturation threshold 0.05, value
threshold 0.95, minimum component size 8, and closing radius 1 for every
hospital. The descriptor summarized tissue fraction, occupancy on a $3\times3$
grid, connected-component count and area, perimeter, compactness, eccentricity,
and centroid. Labels were matched exactly. Empty masks used prespecified
sentinels and a missingness indicator. Matching always used the benchmark
periphery $\Omega\setminus\Omega_R$, so buffer analyses did not change donors.
Mask-stability checks perturbed the thresholds without using classifier
outputs. Following established matching practice, selection used no classifier
outcomes, and downstream interpretation reports balance and common-support
coverage rather than extrapolating beyond accepted triplets
\cite{stuart2010matching}.

For source $i$ in hospital $a$ and target hospital $b$, an eligible triplet
contained a same-hospital donor $j$ and cross-hospital donor $k$ with the same
label and no source-slide overlap. When patient identifiers were recoverable,
the same exclusion rule was applied at patient level. Standardized descriptor
distances were denoted $d_{ij}$ and $d_{ik}$, and triplets were jointly ranked by
\begin{equation}
\mathcal{J}(i,j,k)=d_{ij}+d_{ik}
+\lambda_{\mathrm{bal}}|d_{ij}-d_{ik}|
+\lambda_{\mathrm{pair}}
\|\widetilde{\mathbf v}_j-\widetilde{\mathbf v}_k\|_2.
\label{eq:triplet_objective}
\end{equation}
The frozen matcher selected $K=3$ triplets per direction with
$\lambda_{\mathrm{bal}}=2.0$, $\lambda_{\mathrm{pair}}=0.25$,
$\tau_d=6.0$, $\tau_{\mathrm{bal}}=1.0$, a 64-candidate pool, donor reuse cap
20, and deterministic tie-breaking. Sources without three acceptable pairs
were excluded from both arms. The training hospitals $\{0,3,4\}$ supplied the
cross-donor bank; the six audit directions were
$1\rightarrow\{0,3,4\}$ and $2\rightarrow\{0,3,4\}$. A fixed, stratified
10,000-source subset per held-out hospital was sampled with seed 2026 after
matching eligibility and reused in every construction analysis.

Peripheral tumor extent could not be included because exact alignment between
whole-slide annotations and WILDS patches was not established. We therefore
omit lesion-aware results and treat residual peripheral pathology as an
unmeasured matching factor. Figure~\ref{fig:triplet_matching} summarizes the
matched-triplet selection procedure.

\begin{figure}[H]
\centering
\includegraphics[width=\columnwidth,height=4.1cm,keepaspectratio]{figures/figure4_triplet_matching.png}
\caption{Balanced matched-triplet selection. Same-hospital and cross-hospital
donors are jointly scored under distance, balance, eligibility, and reuse
constraints. Donors must come from a different slide; a different patient is
also required when patient identifiers are available. Classifier outputs are
not used for matching.}
\label{fig:triplet_matching}
\end{figure}

\subsection{Counterfactual Arms and Quality Gates}
\label{subsec:diagnostics}

Each accepted triplet produced the original source, a same-hospital composite,
and a cross-hospital composite. The two composites used identical protection
and blending rules, and the same-hospital arm controlled for generic
replacement and boundary artifacts. A failed donor or construction check
removed the paired arms together. Tensor-level ROI identity was binding; any
violation invalidated the run.

Before prediction analysis, we measured directed coverage as
\begin{equation}
\mathrm{Coverage}_{a\rightarrow b}
=|\mathcal I_{ab}|/|\mathcal I_a^{\mathrm{cand}}|,
\label{eq:pair_coverage}
\end{equation}
where $\mathcal I_a^{\mathrm{cand}}$ is the eligible source set and
$\mathcal I_{ab}$ is the accepted set. We also reported common-support
coverage, restricted to sources matched to every prespecified donor hospital,
and the equal-weight mean of directed coverages. Balance was assessed with the
paired standardized mean difference
\begin{equation}
\operatorname{pSMD}(u)=
\overline{(u_{\mathrm{cross}}-u_{\mathrm{within}})}
/\sigma_{u,\mathrm{train}}.
\end{equation}
The confirmatory gate required directed and common-support coverage of at
least 0.90 and absolute pSMD below 0.10 overall and for every feature. Boundary
quality was checked by comparing cross-seam intensity and gradient
discontinuities and by visually reviewing a prespecified sample.

\subsection{Audit Outcomes}
\label{subsec:metrics}

For each seed, logits were normalized with the median and interquartile range
of ID-validation originals:
\begin{equation}
\widetilde z_\theta(\mathbf{x})=
\frac{z_\theta(\mathbf{x})-\operatorname{median}(z_\theta)}
{\max\{\operatorname{IQR}(z_\theta),0.001\}}.
\end{equation}
For each source and directed pair, $\delta_i^{\mathrm W}$ and
$\delta_i^{\mathrm X}$ were the mean absolute normalized-logit changes across
the $K$ same-hospital and cross-hospital donors. The directed Hospital Context
Sensitivity was
\begin{equation}
\mathrm{HCS}_{a\rightarrow b}=
\frac{1}{N_{ab}}\sum_{i\in\mathcal I_{ab}}
(\delta_i^{\mathrm X}-\delta_i^{\mathrm W}).
\label{eq:pair_hcs}
\end{equation}
Thus, positive HCS indicates greater instability under cross-hospital context;
negative HCS indicates greater instability under matched same-hospital
replacement. Hospital Context Effect (HCE) retained direction by replacing
the two absolute changes with the corresponding signed cross-minus-within
logit shifts. Global summaries were computed either on the common-support
source set or as an equal-weight mean of directed estimates. We never pooled
hospital-1 and hospital-2 source populations or reweighted pairs using observed
outcomes.

\subsection{Statistical Analysis}
\label{subsec:inference}

The source ROI was the paired analysis unit. Reuse of source, same-hospital
donor, and cross-hospital donor identities created crossed dependence, a form
of multiway clustering that invalidates independence-based intervals
\cite{cameron2011multiway}. For
each seed, 95\% intervals therefore used 2,000 product-exponential multiplier
bootstrap draws over physical identities
\cite{owen2007pigeonhole,owen2012bootstrapping}; a reused
identity received the same draw in every role, and triplets were not rematched.
Patient identity was used when recoverable and slide identity otherwise.
Model seed was not bootstrapped because only three seeds were available.

We report all three seed-specific estimates and intervals, followed by the
mean and standard deviation of the three point estimates. Effect size,
intervals, coverage, and diagnostic validity were primary; nominal tests and
Benjamini--Hochberg adjustment were secondary \cite{benjamini1995fdr}.
Figure~\ref{fig:bootstrap_dependencies} illustrates the crossed weighting
structure used by the bootstrap.

\begin{figure}[H]
\centering
\includegraphics[width=\columnwidth,height=3.7cm,keepaspectratio]{figures/figure5_crossed_bootstrap.png}
\caption{Crossed physical-identity bootstrap. Reused source and donor
identities share multiplicity weights. Inference is performed separately for
each model seed.}
\label{fig:bootstrap_dependencies}
\end{figure}

\section{Results}

\subsection{Evaluation Protocol and Analysis Scope}

Following the paired audit design, the evaluation asks whether the matching
procedure supports a valid comparison, whether cross-hospital context changes
the model beyond same-hospital replacement, and whether the observed response
persists across seeds, hospital directions, and construction settings. All
analyses use the same frozen classifiers, source subsets, donor rules, and
metrics described in Section~\ref{sec:methods}. We use a two-stage evaluation: first, coverage,
balance, and ROI-identity checks establish whether an audit population is
admissible; second, global, directed, and robustness analyses characterize the
model response. Hospital-1 validation is confirmatory. Because the frozen
matcher fails the hospital-2 coverage gate, all adapted H2 results are reported
as post-hoc exploratory evidence.

\subsection{Matching Feasibility and Balance}

Table~\ref{tab:matching_balance} reports coverage and within/cross
comparability before any prediction analysis. Acceptance criteria were fixed
from donor metadata and did not depend on hospital-specific outcomes. Detailed
feature balance, reuse distributions, and attrition are provided in the
supplement.

\begin{table}[H]
\caption{Matching feasibility and balance. H1 uses the frozen confirmatory
matcher. The frozen H2 gate failed because common-support coverage was below
0.90; H2 rows report the clearly labeled post-hoc adapted matcher (compactness
caliper 0.10, candidate pool 128). H1 and H2 candidate counts are 34,904 and
85,054; common-support coverage is 97.26\% and 82.10\%, respectively.}
\label{tab:matching_balance}
\centering
\scriptsize
\setlength{\tabcolsep}{2.4pt}
\begin{tabular}{@{}lrrrr@{}}
\toprule
Pair & Accepted & Coverage & Dist. pSMD & Max $|\mathrm{pSMD}|$ \\
\midrule
$1\rightarrow0$ & 34,196 & 97.97\% & $-0.0333$ & 0.0187 \\
$1\rightarrow3$ & 34,009 & 97.44\% & $-0.0386$ & 0.0154 \\
$1\rightarrow4$ & 34,184 & 97.94\% & $-0.0362$ & 0.0140 \\
\midrule
$2\rightarrow0^\dagger$ & 69,964 & 82.26\% & 0.0503 & 0.0556 \\
$2\rightarrow3^\dagger$ & 81,907 & 96.30\% & $-0.0083$ & 0.0138 \\
$2\rightarrow4^\dagger$ & 82,026 & 96.44\% & $-0.0115$ & 0.0095 \\
\bottomrule
\end{tabular}
\end{table}

H1 comfortably passed every frozen gate: common-support coverage was 97.26\%,
directed coverage exceeded 97.4\%, and the largest absolute feature pSMD was
0.0187. The original frozen matcher did not yield acceptable H2 common
support. We therefore do not present a confirmatory H2 estimate. The adapted
H2 matcher improved feasibility but retained lower common-support coverage
(82.10\%), driven by $2\rightarrow0$; all subsequent H2 results are post-hoc
exploratory and cannot be used to validate the original confirmatory claim.

\subsection{Primary Context-Sensitivity Results}

Table~\ref{tab:erm_sensitivity} reports the locked audit for all three
ResNet-50 seeds. The primary evidence comprises the seed-specific effect
sizes and 95\% intervals, the across-seed mean and standard deviation, and the
corresponding matching coverage. Negative estimates are retained rather than
truncated or interpreted as null results.

\begin{table}[H]
\caption{Global ResNet-50 paired context audit over prespecified seeds 11,
42, and 101. Estimates are pair-weighted; common-support estimates were
nearly identical. Mean$\pm$SD summarizes three seed point estimates; detailed
seed values and intervals appear in Table~\ref{tab:seed_results}.}
\label{tab:erm_sensitivity}
\centering
\scriptsize
\setlength{\tabcolsep}{3.0pt}
\resizebox{\columnwidth}{!}{%
\begin{tabular}{@{}lccc@{}}
\toprule
Population / class & $\mathrm{Cov.}_{\cap}/\mathrm{Cov.}_{\mathrm{pair}}$
& HCS mean$\pm$SD & HCE mean$\pm$SD \\
\midrule
H1 / confirmatory & 97.26\% / 97.78\% & $-.0906\pm.0097$ & $.0700\pm.0584$ \\
H2$^\dagger$ / exploratory & 82.10\% / 91.67\% & $.1180\pm.0694$ & $.2394\pm.1556$ \\
\bottomrule
\end{tabular}
}
\end{table}

For confirmatory H1, all three global HCS intervals excluded zero:
$-0.0959$ [$-0.1330,-0.0458$], $-0.0793$
[$-0.1120,-0.0420$], and $-0.0965$ [$-0.1300,-0.0585$]. Thus cross-hospital
transplants induced \emph{less} absolute normalized-logit change than their
matched same-hospital controls. Common-support results differed by less than
0.0007 in every seed, excluding source-set composition as an explanation.
H1 HCE was positive on average but every seed-specific interval included
zero, so the signed direction was not stable. Exploratory H2 HCS was positive
on average, but seed 11 included zero; the H2 HCE interval excluded zero only
for seed 11. These results support heterogeneity, not a single benchmark-wide
hospital-context effect.

\subsection{Hospital-Level Heterogeneity}

We next examine whether the global estimates conceal hospital-specific
variation. The only eligible directions are
$1\rightarrow0$, $1\rightarrow3$, and $1\rightarrow4$ for OOD validation,
and $2\rightarrow0$, $2\rightarrow3$, and $2\rightarrow4$ for final OOD
test. Hospitals $\{0,3,4\}$ never appear as audit sources. Pair-specific HCS,
HCE, and matching coverage are reported in Table~\ref{tab:hospital_pairs};
Fig.~\ref{fig:directed_matrices} shows the seed-specific confidence intervals.
Every directed display is split-specific: hospital-1 OOD-validation sources
and hospital-2 final-test sources are reported separately and never pooled.
Hospital-2 predictions were opened only after pipeline freeze. The frozen
confirmatory matcher failed its coverage gate; the subsequent adapted H2
matcher did not alter the models, metrics, audit subset, or H1 analysis and is
reported only as post-hoc exploratory.

\begin{table}[H]
\caption{Directed HCS/HCE for hospital 1 OOD validation and hospital 2 final
OOD test. Values are across-seed mean$\pm$SD; seed-specific estimates and 95\%
intervals are displayed in Fig.~\ref{fig:directed_matrices}.}
\label{tab:hospital_pairs}
\centering
\scriptsize
\setlength{\tabcolsep}{2.4pt}
\resizebox{\columnwidth}{!}{%
\begin{tabular}{@{}lrrrr@{}}
\toprule
Pair & Accepted & Coverage & HCS & HCE \\
\midrule
$1\to0$ & 34,196 & 97.97\% & $-.1047\pm.0058$ & $.0910\pm.0395$ \\
$1\to3$ & 34,009 & 97.44\% & $-.0631\pm.0053$ & $.1229\pm.0289$ \\
$1\to4$ & 34,184 & 97.94\% & $-.1039\pm.0212$ & $-.0038\pm.1167$ \\
\midrule
$2\to0^\dagger$ & 69,964 & 82.26\% & $.1103\pm.0733$ & $.2616\pm.1736$ \\
$2\to3^\dagger$ & 81,907 & 96.30\% & $.1593\pm.0890$ & $.2972\pm.1533$ \\
$2\to4^\dagger$ & 82,026 & 96.44\% & $.0844\pm.0592$ & $.1594\pm.1678$ \\
\bottomrule
\end{tabular}
}
\end{table}

\begin{figure}[H]
\centering
\includegraphics[width=0.98\textwidth]{figures/directed_effects.pdf}
\caption{Per-seed directed HCS/HCE with 95\% crossed physical-identity
multiplier-bootstrap intervals. Confirmatory H1 HCS is negative for every
seed and hospital pair, with all intervals excluding zero. Exploratory H2 is
positive for seeds 42 and 101 but less stable for seed 11; HCE is generally
less precise. Shading and $\dagger$ mark the post-hoc H2 analysis.}
\label{fig:directed_matrices}
\end{figure}

\subsection{Robustness and Implementation Checks}

We compare the primary Balanced Matched-Triplet audit with Random Paired
selection, an ROI-only implementation falsification, the same triplets with
an 8-pixel protected source buffer, and hard versus primary feathered
compositing. Random Paired controls generic within/cross replacement but
does not explicitly balance measured morphology. Balanced Matched-Triplet
selection additionally balances tissue geometry and within/cross matching
difficulty. Every row reports its corresponding coverage; construction
variants use the same fixed audit subset and triplets.

\begin{table}[H]
\caption{Matching and construction robustness. HCS/HCE are across-seed
mean$\pm$SD of pair-weighted estimates. Construction variants reuse the
fixed 10,000-source audit subset and matched triplets. H2 is post-hoc
exploratory.}
\label{tab:identification_summary}
\centering
\scriptsize
\setlength{\tabcolsep}{2.2pt}
\begin{tabularx}{\columnwidth}{@{}lXcc@{}}
\toprule
Population / setting & Specification & HCS & HCE \\
\midrule
H1 Balanced & $r=0$, feathered & $-.0906\pm.0097$ & $.0700\pm.0584$ \\
H1 Random Paired & $r=0$, feathered & $-.0780\pm.0247$ & $.0243\pm.0796$ \\
H1 ROI-only & fixed input periphery & $0.0000\pm.0000$ & $0.0000\pm.0000$ \\
H1 protected buffer & $r=8$, feathered & $-.1019\pm.0063$ & $.0434\pm.0480$ \\
H1 hard seam & $r=0$, hard & $-.0866\pm.0076$ & $.0704\pm.0548$ \\
\midrule
H2$^\dagger$ Balanced & $r=0$, feathered & $.1180\pm.0694$ & $.2394\pm.1556$ \\
H2$^\dagger$ Random Paired & $r=0$, feathered & $.0742\pm.0959$ & $.2287\pm.1432$ \\
H2$^\dagger$ ROI-only & fixed input periphery & $0.0000\pm.0000$ & $0.0000\pm.0000$ \\
H2$^\dagger$ protected buffer & $r=8$, feathered & $.1073\pm.0452$ & $.2125\pm.1231$ \\
H2$^\dagger$ hard seam & $r=0$, hard & $.1220\pm.0717$ & $.2370\pm.1539$ \\
\bottomrule
\end{tabularx}
\end{table}

The confirmatory H1 sign and magnitude persisted at $r=8$ and under a hard
boundary, so the negative HCS is not restricted to the immediate ROI boundary
or one seam construction. Random pairing also remained negative but was more
seed-variable, indicating that explicit balancing improves precision more
than it changes direction. The ROI-only control returned exactly zero HCS and
HCE for every seed and both populations, providing a strong implementation
falsification. Exploratory H2 remained positive on average for the buffer and
hard-boundary variants, with substantial seed variability. None of these
checks identifies a spatial or visual causal mechanism.

\section{Discussion}

\subsection{Scope of Interpretation}

The study measures whether a trained tumor classifier changes its prediction
when measured-content-balanced peripheral context is transplanted across
hospital-indexed donor pools while the label-defining central region remains
fixed. Positive HCS indicates more instability under matched cross-hospital
than matched same-hospital transplantation; negative HCS indicates the
reverse. HCE retains the signed paired shift. Neither metric identifies which
scanner, stain, workflow, or biological factor generated the difference.

The central confirmatory result is negative, not null: for H1, every seed and
directed donor hospital produced HCS below zero with a 95\% interval excluding
zero. The model is sensitive to context replacement, but a comparably matched
same-hospital replacement is consistently more destabilizing than the
cross-hospital transplant. This directly falsifies the simple expectation
that cross-hospital peripheral context should uniformly amplify prediction
instability. It does not prove that the model ignores hospital-associated
signals; rather, it shows that the measured response depends on the paired
intervention and source population.

The H2 analysis reinforces the need for that restraint. After the frozen
confirmatory matching gate failed, an adapted post-hoc matcher produced a
positive average HCS, but with lower common-support coverage and visible seed
heterogeneity. H2 therefore generates a hypothesis about population-specific
context response; it cannot rescue or overturn the confirmatory H1 finding.

Common-support summaries hold the source set fixed across donor hospitals
but apply only to jointly matchable sources. Pair-weighted summaries retain
more eligible comparisons while permitting pair-specific source composition.
Coverage and directed estimates therefore accompany both. The Random Paired
comparison indicates how much explicit morphology and difficulty balancing
changes HCS/HCE, without converting the audit into a causal analysis.

The ROI-only check validates implementation, while the protected-buffer and
boundary comparisons test two limited construction alternatives. These
checks can strengthen a scoped context-sensitivity interpretation, but they
cannot establish that a specific acquisition factor caused the response or
that the model uses every form of medical shortcut. Camelyon17-WILDS is
particularly suitable because it provides an explicit label-defining ROI;
diffuse or uncertain targets require another preservation rule.

\subsection{Limitations}

The audit requires a known label-defining ROI, and peripheral transplantation
is not equivalent to rescanning the same patient. Matching controls measured
tissue geometry but cannot remove unmeasured pathology or identify a causal
scanner, stain, or workflow mechanism. Exact alignment of peripheral lesion
annotations was unavailable, and compositing itself remains an artificial
intervention.

Generality is also limited by one benchmark, one convolutional architecture,
three training seeds, and a compute-limited audit subset drawn from only 50
whole-slide images. The frozen H2 matching gate failed; the adapted H2 analysis
is therefore post-hoc and has only 82.10\% common-support coverage.

\section{Conclusions}

MedStyleAudit isolates peripheral-context sensitivity while preserving the
label-defining ROI exactly and controlling for generic replacement effects.
Confirmatory H1 HCS was consistently negative
($-0.0906\pm0.0097$), persisted under buffer and boundary changes, and was
accompanied by exact-zero ROI-only controls. Post-hoc H2 estimates were positive
but heterogeneous. The results therefore reject a uniform hospital-context
shortcut signature and support reporting matched controls, coverage gates,
and population-specific effects. MedStyleAudit remains a sensitivity audit,
not a method for identifying causal acquisition mechanisms.

%%%%%%%%%%%%%%%%%%%%%%%%%%%%%%%%%%%%%%%%%%
\vspace{6pt}

% Complete the contribution and funding fields before submission.
\authorcontributions{To be completed by the authors according to the CRediT taxonomy.}
\funding{To be completed by the authors.}
\institutionalreview{The original CAMELYON collection was approved by the Commissie Mensgebonden Onderzoek regio Arnhem--Nijmegen (protocol 2016-2761), as reported by Litjens et al. The present study used only publicly released, de-identified Camelyon17-WILDS data and involved no new participant recruitment, specimen collection, or access to direct identifiers.}
\informedconsent{Patient consent for the original CAMELYON collection was waived by the approving ethics committee, as reported by Litjens et al. No participants were recruited for the present secondary analysis.}
\dataavailability{The Camelyon17-WILDS benchmark is available through WILDS at \url{https://wilds.stanford.edu/datasets/} under a CC0 public-domain dedication. Code is available at \url{https://github.com/PeiYuanHao/MedStyleAudit-github}; experimental artifacts are available at \url{https://huggingface.co/datasets/PeiyuanHao/MedStyleAudit-Experiments}.}
\acknowledgments{During manuscript preparation, the authors used OpenAI Codex for language editing and structural revision. The authors reviewed and edited the output and take full responsibility for the content of the publication.}
\conflictsofinterest{To be completed by the authors.}

\abbreviations{Abbreviations}{%
The following abbreviations are used in this manuscript:\\

\noindent
\begin{tabular}{@{}ll}
HCS & Hospital Context Sensitivity\\
HCE & Hospital Context Effect\\
ROI & Region of interest\\
OOD & Out of distribution\\
AUROC & Area under the receiver operating characteristic curve\\
\end{tabular}
}

\appendixtitles{yes}
\appendixstart
\appendix


\section{Supplementary Protocol and Quality-Assurance Detail}

\subsection{Inference Implementation}

For each model seed separately, independent multiplicity weights are drawn
for every physical source, within-donor, and cross-donor identity within
hospital, and their product weights the paired triplet record. If a physical
identity appears in multiple roles, the same draw is reused for that
identity; if paired outcomes are stored in two arm rows, both rows receive
the same joint weight. Model seed is not a bootstrap cluster. The three
seed-specific estimates and intervals are reported individually, followed
by the mean and standard deviation of the three point estimates. We use 2,000
product-exponential multiplier-bootstrap draws and the 2.5th and 97.5th
percentiles for each 95\% interval, as frozen in
\texttt{configs/final/FINAL\_PROTOCOL.yaml}.
Complete nominal and Benjamini--Hochberg-adjusted test results, if computed,
remain secondary to effect sizes and intervals.

\subsection{Protocol Tables}

\begin{table}[H]
\caption{Prespecified descriptor components. ``Optional'' requires exact
lesion-coordinate validation.}
\label{tab:feature_membership}
\centering
\scriptsize
\setlength{\tabcolsep}{2.3pt}
\begin{tabularx}{\columnwidth}{@{}Xc@{}}
\toprule
Feature & Primary content match \\
\midrule
Label (exact constraint) & Yes \\
Tissue fraction / spatial occupancy & Yes \\
Component areas, count, perimeter & Yes \\
Compactness / eccentricity / centroid & Yes \\
Peripheral tumor fraction / presence & Optional \\
\bottomrule
\end{tabularx}
\end{table}

\begin{table}[H]
\caption{Canonical frozen final-protocol fields synchronized with
\texttt{configs/final/FINAL\_PROTOCOL.yaml}.}
\label{tab:full_implementation}
\centering
\scriptsize
\setlength{\tabcolsep}{2.5pt}
\begin{tabularx}{\columnwidth}{@{}>{\raggedright\arraybackslash}Xl@{}}
\toprule
Configuration & Value \\
\midrule
Patch / label ROI & $96\times96$ / central $32\times32$ \\
Architecture & ResNet-50 \\
Initialization / loss & random / BCE \\
Audit sources & hospital 1 OOD validation; hospital 2 final OOD test \\
Cross-donor bank & training hospitals $\{0,3,4\}$ \\
Mask filter / thresholds & median 3; $S\geq.05$; $V\leq.95$; min object 8; closing radius 1 \\
Optimizer / schedule & AdamW / cosine \\
Learning rate / weight decay & $3\times10^{-4}$ / $10^{-4}$ \\
Batch size / epochs & 128 / 20 \\
Donors per ROI, $K$ & 3 \\
Distance / balance thresholds & $\tau_d=6.0$, $\tau_{\mathrm{bal}}=1.0$ \\
Triplet weights / calipers & $\lambda_{\mathrm{bal}}=2.0$, $\lambda_{\mathrm{pair}}=.25$; none primary \\
Candidate pool & 64 primary; 128 post-hoc H2$^\dagger$ \\
Donor / donor-slide reuse caps & 20 / 200,000 \\
Tie-breaking & cost, within ID, cross ID \\
Protection margin, $r$ & primary 0; robustness 8 \\
Normalized-logit floor, $q_{\min}$ & 0.001 \\
Feather width, $w$ & 4 \\
Boundary comparison & hard vs. primary feathered \\
Random seeds & 11, 42, 101 \\
Bootstrap draws (within each seed) & 2,000 \\
Bootstrap identities & physical source / within donor / cross donor \\
Bootstrap model seed as cluster & No \\
Interval type & percentile, crossed multiplier bootstrap \\
Across-seed summary & three estimates; mean $\pm$ standard deviation \\
Model-selection criterion & H1 OOD-validation AUROC \\
Final-test evaluation status & complete; H2$^\dagger$ exploratory \\
\bottomrule
\end{tabularx}
\end{table}

\begin{table}[H]
\caption{ROI-preservation and null-control checks. Exact ROI equality at the
final model input was a binding requirement; any violation would terminate the
run. All 30 configuration--seed evaluations completed without a violation. In
the ROI-only control, donor periphery is excluded, so a null HCS/HCE response is
the expected invariant output rather than an estimated biological effect.}
\label{tab:preservation}
\centering
\scriptsize
\setlength{\tabcolsep}{2.8pt}
\resizebox{\columnwidth}{!}{%
\begin{tabular}{@{}lcc@{}}
\toprule
Check & H1 confirmatory & H2$^\dagger$ exploratory \\
\midrule
Final-input ROI identity & Passed (15/15) & Passed (15/15) \\
ROI-only null control (HCS/HCE) & Confirmed (3/3 seeds) &
Confirmed (3/3 seeds) \\
\bottomrule
\end{tabular}
}
\end{table}

\begin{table}[H]
\caption{Detailed matching attrition and balance. Max reuse is the frozen
per-donor cap; H2 uses the post-hoc adapted matcher.}
\label{tab:supp_matching}
\centering
\scriptsize
\begin{tabular}{@{}lrrrrr@{}}
\toprule
Pair & Cand. & Accept. & Exclusion & Cap & Max $|\mathrm{pSMD}|$ \\
\midrule
$1\rightarrow0$ & 34,904 & 34,196 & gate & 20 & .0187 \\
$1\rightarrow3$ & 34,904 & 34,009 & gate & 20 & .0154 \\
$1\rightarrow4$ & 34,904 & 34,184 & gate & 20 & .0140 \\
$2\rightarrow0^\dagger$ & 85,054 & 69,964 & gate & 20 & .0556 \\
$2\rightarrow3^\dagger$ & 85,054 & 81,907 & gate & 20 & .0138 \\
$2\rightarrow4^\dagger$ & 85,054 & 82,026 & gate & 20 & .0095 \\
\bottomrule
\end{tabular}
\end{table}

\begin{table}[H]
\caption{Individual ResNet-50 seed results. HCS/HCE are pair-weighted global
estimates with seed-specific 95\% crossed multiplier-bootstrap intervals.
The AUROC column reports the H1 OOD-validation model-selection score and is
therefore repeated for the H2 rows. H2 is post-hoc exploratory.}
\label{tab:seed_results}
\centering
\scriptsize
\resizebox{\columnwidth}{!}{%
\begin{tabular}{@{}llccc@{}}
\toprule
Population & Seed & H1 OOD-val AUROC & HCS [95\% CI] & HCE [95\% CI] \\
\midrule
H1 & 11 & .8688 & $-.0959\ [-.1330,-.0458]$ & $.0986\ [-.1075,.2603]$ \\
H1 & 42 & .8832 & $-.0793\ [-.1120,-.0420]$ & $.1086\ [-.0859,.2572]$ \\
H1 & 101 & .8983 & $-.0965\ [-.1300,-.0585]$ & $.0028\ [-.1747,.1701]$ \\
H2$^\dagger$ & 11 & .8688 & $.0514\ [-.1111,.1743]$ & $.4061\ [.1526,.5708]$ \\
H2$^\dagger$ & 42 & .8832 & $.1127\ [.0230,.1910]$ & $.0980\ [-.1706,.3029]$ \\
H2$^\dagger$ & 101 & .8983 & $.1900\ [.0613,.3155]$ & $.2141\ [-.0751,.4629]$ \\
\bottomrule
\end{tabular}
}
\end{table}

\FloatBarrier

\FloatBarrier
\begin{adjustwidth}{-\extralength}{0cm}
\reftitle{References}

\bibliography{references}

\end{adjustwidth}

\PublishersNote{}
\end{document}
